% -----------------------------------------------------------
\section{Comprehend}
% -----------------------------------------------------------
	\label{COMPREHEND}
	
% % ----------------------
% \subsection{See also}
% % ----------------------

% ----------------------
\subsection{1828 American Dictionary of the English Language} % *1828
% ----------------------
	\label{COMPREHEND-1828}

	\begin{compactitem}
		\item[1] To contain; to include; to comprise.
		\item[2] To imply; to contain or include by implication or construction.
		\item[3] To understand; to conceive; that is, to take, hold or contain in the mind; to possess or to have in idea.
	\end{compactitem}

% % ----------------------
% \subsection{Oxford English Dictionary} % *OED
% % ----------------------
	% \label{COMPREHEND-OED}

	% \begin{compactitem}
		% \item[]
	% \end{compactitem}

% ----------------------
\subsection{Scriptural Usage} % *SCRIPTURES
% ----------------------
	\label{COMPREHEND-SCRIPTURES}

	% -----
	\subsubsection{Old Testament} % *OT
	% -----
		\label{COMPREHEND-OT}
	
		% --
		\paragraph{KJV}
		% --
			\label{COMPREHEND-OT-KJV}
	
			\begin{tabular}{ll}
				Total occurrences in the OT & 2
			\end{tabular}
			
			\begin{compactdesc}
				\item[{\hyperref[JOB-37-5]{Job 37:5}}]
				\item[{\hyperref[ISAIAH-40-12]{Isaiah 40:12}}]
			\end{compactdesc}
			
		% % --
		% \paragraph{Hebrew Bible}
		% % --
			% \label{COMPREHEND-HEBREW}
		
			% %
			% \subparagraph{\texorpdfstring{\he{sample word}}{}}
			% %
				% \label{PAGA} % whatever is in step bible without periods
			
				% \begin{compactdesc}
					% \item[HALOT , PDF ] \textbf{}
						% \begin{compactitem}
							% \item[1]
						% \end{compactitem}
					% \item[BDB , PDF ] \textbf{}
						% \begin{compactitem}
							% \item[1]
						% \end{compactitem}
					% \item[DCH PDF ] \textbf{}
						% \begin{compactitem}
							% \item[1]
						% \end{compactitem}
				% \end{compactdesc}
				
				% \begin{compactdesc}
					% \item[Some OT uses] \textbf{}
						% \begin{compactdesc}
							% \item[]
						% \end{compactdesc}
				% \end{compactdesc}
			
		% % --
		% \paragraph{Septuagint}
		% % --
			% \label{COMPREHEND-LXX}
		
			% %
			% \subparagraph{\texorpdfstring{\gr{sample word}}{}}
			% %
		
	% -----
	\subsubsection{New Testament} % *NT
	% -----
		\label{COMPREHEND-NT}
	
		% --
		\paragraph{KJV}
		% --
			\label{COMPREHEND-NT-KJV}
			
	
			\begin{tabular}{ll}
				Total occurrences in the NT & 3
			\end{tabular}
			
			\begin{compactdesc}
				\item[{\hyperref[JOHN-1-5]{John 1:5}}]
				\item[{\hyperref[ROMANS-13-9]{Romans 13:9}}]
					\begin{compactitem}
						\item The Greek word here is \gr{ἀνακεφαλαιόομαι}, which just means ``to summarize.''
					\end{compactitem}
				\item[{\hyperref[EPHESIANS-3-18]{Ephesians 3:18}}]
			\end{compactdesc}
			
		% --
		\paragraph{Greek New Testament} % *GREEK
		% --
			\label{COMPREHEND-GREEK}
		
			%
			\subparagraph{\texorpdfstring{\gr{καταλαμβάνω}}{}}
			%
				\label{KATALAMBANO}
			
				\begin{compactdesc}
					\item[BDAG 460b] \textbf{}
						\begin{compactitem}
							\item[1] to make something one's own, \textit{win, attain} \textit{(There is a lengthy discussion here of the meaning in John 1:5, starting at the beginning of page 520, PDF 301; ``the darkness did not grasp it,'' ``the darkness did not overcome it,'' though I definitely like ``attain'' and ``win'' here too.)}
							\item[2] to gain control of someone through pursuit, \textit{catch up with, seize}
							\item[3] to come upon someone, with implication of surprise, \textit{catch}
							\item[4] to process information, \textit{understand, grasp}
								\begin{compactitem}
									\item John 1:5 is mentioned here but it refers you back to meanings \#1 and \#2.
								\end{compactitem}
						\end{compactitem}
					% \item[LSJ , PDF ] \textbf{}
						% \begin{compactitem}
							% \item 
						% \end{compactitem}
					\item[NIDNTTE 3:83] ``\gr{καταλαμβάνω} occurs 13x (plus 2x in the story of the woman `caught' in adultery, \hyperref[JOHN-8-3]{John 8:3--4}). In negative contexts, it is used of an evil spirit ``seizing'' a boy and throwing him to the ground (\hyperref[MARK-9-18]{Mark 9:18}), of darkness not `overcoming' or `understanding' the light (\hyperref[JOHN-1-5]{John 1:5}; similar to \hyperref[JOHN-12-35]{John 12:35}), and of the final day `overtaking' unprepared Christians (`surprise you like a thief,' \hyperref[1THESSALONIANS-5-4]{1 Thessalonians 5:4}). The remaining occurrences are positive. Paul does not claim that he has `already obtained [\gr{λαμβάνω}]' or `taken hold of [\gr{καταλαμβάνω}]' the goal of his high calling (\hyperref[PHILIPPIANS-3-12]{Philippians 3:12a, 13}, but rather: `I press on to take hold of that for which Christ Jesus took hold of me [\gr{διώκω δὲ εἰ καὶ καταλάβω, ἐφ᾽ ᾧ καὶ κατελήμφθην ὑπὸ Χριστοῦ}]' (3:12b). In contrast to Israelites who have not responded to the gospel, Gentiles have `obtained' righteousness through faith (\hyperref[ROMANS-9-30]{Romans 9:30}; NRSV, `attained'). Believers are exhorted to run their spiritual race so they can `win' the victor's imperishable crown of eternal life (\hyperref[1CORINTHIANS-9-24]{1 Corinthians 9:24}, in parallel with \gr{λαμβάνω}). The middle voice can denote mental grasping, thus `to comprehend, understand, realize' (\hyperref[ACTS-4-13]{Acts 4:13}; \hyperref[ACTS-10-34]{10:34}; \hyperref[ACTS-25-25]{25:25}), and it is used in \hyperref[EPHESIANS-3-18]{Ephesians 3:18} of understanding the extent of God's love (cf. English \textit{comprehend}, from Latin \textit{prehendere}, `to grasp').''
					\item[TDNT 4:12] \phantom{.}
						\begin{compactitem}
							\item ``\gr{κατα}--- orig. `from above to below,' hence completely, so that \gr{καταλαµβάνω} is a strengthening of the simple form.
								\begin{compactitem}
									\item Note that the use of ``comprehend'' in \hyperref[DC88-6]{DC 88:6} involves Christ ``ascending on high'' and ``descending below'' all things, in keeping with the effect that \gr{κατα} has on \gr{λαµβάνω}, if we read that verse as referring to \hyperref[JOHN-1-5]{John 1:5}.
								\end{compactitem}
							\item ``In the NT a point which has emerged already from this review comes out with particular clarity, namely, that \gr{κατά} gives to the simple form the character either of intensity (to grasp with force, Mt. 9:18) or of suddenness (to surprise, 1 Th. 5:4 [Jn. 8:3 f. R D etc.; 6:17 \he{'} D]; Jn. 12:35 of the time after the death of Jesus). The epistemological nature of the term ($\rightarrow$ c. and I, 692) is not to be seen except in Lk. (Ac.).''
							\item ``The act is used a. in the positive sense of `to attain definitively,' in R. 9:30 in respect of the righteousness of faith, and paradoxically, without any effort (the element of lighting upon or overtaking is still discernible, as in the next passage), or in Phil. 3:12b, 13 in respect of close fellowship with Christ provisionally demonstrated in dying with Him and consummated by the resurrection (with the pneumatic body). The Christian must constantly seek this. He must strive to overtake it. He will finally possess it in the \gr{ἐξανάστασις}. It must be affirmed afresh each day in the life of faith until it is finally consummated in the resurrection. Cf. also Phil. 3:12c, of the Christian whom Christ draws fully into this fellowship of death and resurrection, and 1 C. 9:24 of the Christian's crown of victory, i.e., the consummation of this fellowship == \gr{ζωὴ αἰώνιος}. It is also used in the negative sense `to overpower,' Jn. 1:5: The darkness of separation from God has not succeeded in overcoming the light, the new religious life, which is present in the Logos, in the divine Christ. It has not been able to vanquish the power of His light. By the very existence of this light the whole sphere of night is overcome and deprived of its power.''
							\item ``Corresponding is the use of the middle (only in the intellectual sphere): `to establish' (Ac. 4:13; 25:25), `to grasp fundamentally, `to appropriate to oneself inwardly' (Ac. 10:34; Eph. 3:18: the all-permeating greatness of the \gr{>ag'aph toῦ qristoῦ}.''
						\end{compactitem}
				\end{compactdesc}
				
				\begin{compactdesc}
					\item[Some NT uses] \textbf{}
						\begin{compactdesc}
							\item[John 1:5]
							\item[Ephesians 3:18]
						\end{compactdesc}
				\end{compactdesc}
		
	% -----
	\subsubsection{Book of Mormon} % *BOM
	% -----
		\label{COMPREHEND-BOM}
	
		\begin{tabular}{ll}
			Total occurrences in the BOM & 4
		\end{tabular}
		
		\begin{compactdesc}
			\item[{\hyperref[MOSIAH-4-9]{Mosiah 4:9}}] \textbf{(2)}
			\item[{\hyperref[ALMA-26-35]{Alma 26:35}}]
			\item[{\hyperref[MORMON-9-16]{Mormon 9:16}}]
		\end{compactdesc}
		
	% -----
	\subsubsection{Doctrine and Covenants} % *DC
	% -----
		\label{COMPREHEND-DC}
	
		\begin{tabular}{ll}
			Total occurrences in the DC & 12
		\end{tabular}
		
		\begin{compactdesc}
			\item[{\hyperref[DC6-21]{DC 6:21}}]
			\item[{\hyperref[DC10-58]{DC 10:58}}]
			\item[{\hyperref[DC34-2]{DC 34:2}}]
			\item[{\hyperref[DC39-2]{DC 39:2}}]
			\item[{\hyperref[DC45-7]{DC 45:7}}]
			\item[{\hyperref[DC88-6]{DC 88:6}}]
			\item[{\hyperref[DC88-41]{DC 88:41}}]
			\item[{\hyperref[DC88-43]{DC 88:43}}]
			\item[{\hyperref[DC88-48]{DC 88:48}}]
			\item[{\hyperref[DC88-49]{DC 88:49}}] \textbf{(2)}
			\item[{\hyperref[DC88-67]{DC 88:67}}]
		\end{compactdesc}
		
	% -----
	\subsubsection{Pearl of Great Price} % *PGP
	% -----
		\label{COMPREHEND-PGP}
	
		\begin{tabular}{ll}
			Total occurrences in the PGP & 2
		\end{tabular}
		
		Note that the second use in the PGP is to a JS-H note.
		
		\begin{compactdesc}
			\item[{\hyperref[ABRAHAM-4-4]{Abraham 4:4}}]
		\end{compactdesc}
		
% % ----------------------
% \subsection{Key Phrases} % *KEYPHRASES *PHRASES
% % ----------------------
	% \label{COMPREHEND-PHRASES}

	% \begin{compactdesc}
		% \item[] \textbf{#times} | list
			% % \begin{compactdesc}
				% % \item[Bible anchor verses] \phantom{.}
					% % \begin{compactdesc}
						% % \item[Romans 5:5] \gr{}
					% % \end{compactdesc}
				% % \item[Commentary] ---
			% % \end{compactdesc}
	% \end{compactdesc}
	
% % ----------------------
% \subsection{Bible Dictionaries} % *DICTIONARIES
% % ----------------------
	% \label{COMPREHEND-BD}

% % ----------------------
% \subsection{Restoration Usage} % *RESTORATION
% % ----------------------
	% \label{COMPREHEND-RESTORATION}

	% % -----
	% \subsubsection{Joseph Smith} % *JS
	% % -----
		% \label{COMPREHEND-JS}
	
		% \begin{compactdesc}
			% \item[{\hyperref[KEY]{Date/Source}}]
		% \end{compactdesc}
	
	% % -----
	% \subsubsection{Brigham Young} % *BY
	% % -----
		% \label{COMPREHEND-BY}
	
		% \begin{compactdesc}
			% \item[{\hyperref[KEY]{Date/Source}}]
		% \end{compactdesc}
	
% % ----------------------
% \subsection{Commentary and Studies} % *COMMENTARY *STUDIES
% % ----------------------
	% \label{COMPREHEND-COMMENTARY}

	% % -----
	% \subsubsection{Key Points}
	% % -----
		% \label{COMPREHEND-KEYS}
	
		% \begin{compactitem}
			% \item
		% \end{compactitem}
		
	% % -----
	% \subsubsection{Sample Study}
	% % -----
		% \label{COMPREHEND-study}